\begin{table}[H]
\caption{\label{tab:Summary_task_intensities}Task Intensity of Major Occupations
Categories}

\begin{centering}
\begin{tabular}{l>{\centering}p{1cm}>{\centering}p{2cm}>{\centering}p{2cm}>{\centering}p{2cm}}
\toprule 
 & \emph{RTI }index & Abstract tasks & Routine tasks & Manual tasks\tabularnewline
\cmidrule{2-5} \cmidrule{3-5} \cmidrule{4-5} \cmidrule{5-5} 
Managers/Legislators & - & \cellcolor[gray]{0.75}+ & - & -\tabularnewline
Professionals & - & \cellcolor[gray]{0.75}+ & - & -\tabularnewline
Technicians & - & \cellcolor[gray]{0.75}+ & - & -\tabularnewline
Clerks & + & - & \cellcolor[gray]{0.75}+ & -\tabularnewline
Service and sales & + & \cellcolor[gray]{0.75}+ & - & -\tabularnewline
Craft and related trade & - & - & + & \cellcolor[gray]{0.75}+\tabularnewline
Machine operators & + & - & \cellcolor[gray]{0.75}+ & +\tabularnewline
Elementary occupations & + & - & + & \cellcolor[gray]{0.75}+\tabularnewline
\bottomrule
\end{tabular}
\par\end{centering}
\medskip

\textbf{Notes:}The table constructed by following the work of \citet{2013Autor}.
The plus sign indicates that the average value of particular task
for that major occupation category is greater than the average. Minus
sings indicates that the opposite. Gray shaded boxes are signs of
that particular occupations' dominant task intensity.
\end{table}
